\subsection{从李代数到李群的过渡}

我们将 Hausdorff 级数记为 H(X, Y)（第 II 章 § 6，第4号，定义1）。

引理 2. 设 \(\mathbf{L}\) 是定义在 \(\mathbf{R}\) 或 \(\mathbf{C}\) 上的完备赋范李代数。设 \(\mathbf{G}\) 是 \(x \in \mathbf{L}\) 的集合，其中 \(\| x \| < \frac{1}{3} \log \frac{3}{2}\)。设 \(\theta\) 是映射 \(x \mapsto -x\)，从 \(\mathbf{G}\) 到 \(\mathbf{G}\)。设 \(\mathbf{H}\) 是限制在 \(\mathbf{G} \times \mathbf{G}\) 上的 \(\mathbf{L}\) 的 Hausdorff 函数（第 II 章 § 7，第2号）。

\begin{enumerate}
\def\labelenumi{(\roman{enumi})}
\item
  (G, 0, θ, H) 是李群芽。
\item
  设 \(\phi\) 是从 G 到 L 的恒等映射。\(\phi\) 在 0 处的微分是从可赋范李代数 L(G) 到 L 的同构。
\end{enumerate}
(i) 由第 II 章 § 7 第2号推出。

由于 \(\phi\) 是 G 上的图，在 0 处的微分 \(\psi\) 是 \(\phi\) 给出的可赋范空间同构。另一方面，映射 H 的积分级数展开 \(\mathrm{H} = \sum_{i,j\geqslant 0}\mathrm{H}_{ij}\) 满足 \(\mathrm{H}_{11}(x,y) = \frac{1}{2} [x,y]\)。根据§ 3 命题24，对任意 \(a,b\)，它们属于 \(\mathrm{L}(\mathrm{G})\)，

\[
\psi ([ a, b ]) = \mathrm{H} _ {1 1} (\psi (a), \psi (b)) - \mathrm{H} _ {1 1} (\psi (b), \psi (a)) = [ \psi (a), \psi (b) ]
\]

这证明了 (ii)。

称 G 为由 L 定义的李群芽。

假设 K 是超度量的。设 p 是 K 的剩余域的特征。若 \(p \neq 0\)，令 \(\lambda = |p|^{1/(p-1)}\)；若 p = 0，令 \(\lambda = 1\)。

引理 3. 设 L 是定义在 K 上的完备赋范李代数。设 G 是 \(x \in \mathbf{L}\) 的集合，满足 \(\| x \| < \lambda\)。设 H: G × G → G 是 L 的 Hausdorff 函数（第 II 章 § 8，第3号）。

\begin{enumerate}
\def\labelenumi{(\roman{enumi})}
\item
  在复合律 H 下，G 是一个李群，其中 0 是单位元，对所有 \(x \in G\)，-x 是 x 的逆元。
\item
  设 \(\phi\) 是从 \(\mathbf{G}\) 到 \(\mathbf{L}\) 的恒等映射。\(\phi\) 在 0 处的微分是从可赋范李代数 \(\mathrm{L}(\mathrm{G})\) 到 \(\mathbf{L}\) 的同构。
\item
  对任意 \(\mu \in \mathbf{R}_+^*\)，设 \(\mathrm{G}_{\mu}\) 是 \(x \in \mathbf{L}\) 的集合，满足 \(\| x \| < \mu\)。则这些 \(\mathrm{G}_{\mu}\) 在 \(\mu < \lambda\) 时构成 0 的开闭邻域基本系，并且是 \(\mathbf{G}\) 的子群。
\end{enumerate}

断言 (i) 和 (iii) 由第 II 章 § 8 第3号命题3推出，而 (ii) 可像引理2中那样证明。

称 G 为由 L 定义的李群。

定理 2. 设 \(\mathbf{L}\) 是完备可赋范李代数。存在一个李群芽 \(G\)，使得 \(\mathrm{L}(\mathbf{G})\) 与 \(\mathbf{L}\) 同构。任意两个这样的李群芽局部同构。

第一个断言由引理2和引理3推出。第二个断言由第1号定理1的推论1推出。

推论 1. 设 \(\mathbf{G}\) 是李群。存在一个 \(e\) 的邻域，其中不含异于 \(\{e\}\) 的有限子群。若 \(\mathbf{K} = \mathbf{R}\) 或 \(\mathbf{C}\)，则存在一个 \(e\) 的开邻域，其中不含异于 \(\{e\}\) 的子群。

记 \(\mathrm{L}(G) = \mathrm{L}\)。在 \(\mathbf{L}\) 上取一个定义 \(\mathbf{L}\) 拓扑的范数，使得 \(\| [x,y]\| \leqslant \| x\| \| y\|\) 对所有 \(x,y\) 属于 \(\mathbf{L}\) 均成立。

假设 \(\mathbf{K} = \mathbf{R}\) 或 \(\mathbf{C}\)。设 \(G'\) 是由 L 定义的李群芽。存在一个以 0 为中心的开球 \(U'\)，位于 \(G'\) 中，以及一个同构 \(\phi\)，其定义域是李群芽 \(U'\)，陪域为开邻域 \(U\)，其中包含 \(e\)，且位于 G 中。令 \(V' = \frac{1}{2} U'\)、\(V = \phi(V')\)，设 H 是包含于 V 的 G 的子群，且 \(h \in H\)。令 \(x = \phi^{-1}(h) \in V'\)。若 \(x \neq 0\)，则存在整数 \(n > 0\)，使得 \(x, 2x, \ldots, nx\) 属于 \(V'\)，\((n + 1)x \in U'\)，而 \((n + 1)x \notin V'\)。于是 \(h, h^2, \ldots, h^n\) 属于 V，

\[
h ^ {n + 1} \in \mathrm{U}, \quad h ^ {n + 1} \notin \mathrm{V},
\]

这不可能。因此 \(\mathbf{H} = \{e\}\)。

假设 K 是超度量的。只需在 G 是与 L 相关的李群时证明此推论。若 \(g \in G\)，则 g 在 G 中计算的幂，是 Zg 在 L 中计算的元素。当 \(g \neq e\) 时这些元素彼此不同。因此 G 不含异于 \(\{e\}\) 的有限子群。

推论 2. 设 \(k\) 是 \(\mathbf{K}\) 的非离散闭子域，\(\mathbf{G}\) 是定义在 \(k\) 上的李群，且 \(\mathbf{L} = \mathbf{L}(\mathbf{G})\)。假设 \(\mathbf{L}\) 具有可赋范李 \(\mathbf{K}\)-代数结构 \(\mathbf{L}'\)，它与可赋范李 \(k\)-代数结构相容，并且在 \(\mathbf{G}\) 的伴随表示下不变。则 \(\mathbf{G}\) 上存在唯一一个李 \(\mathbf{K}\)-群结构，它与李 \(k\)-群结构相容，且其李代数为 \(\mathbf{L}'\)。

存在定义在 K 上的李群芽 \(G_{1}\)，使得 \(\mathrm{L}(G_{1}) = \mathrm{L}'\)（定理2）。根据第1号定理1的推论1，将 G 和 \(G_{1}\) 看作李 k-群芽时，它们局部同构。因此，G 中存在 e 的一个开邻域 \(G'\)，以及 \(G'\) 上的李 K-群芽结构，其李代数为 L，并且与定义在 k 上的李群芽结构相容。设 V 是 G 中 e 的对称开邻域，满足 \(V^{2} \subset G'\)。设 \(g \in G\)。则 \(\phi = Int g\) 是从 \(G'\) 的某个充分小的开李子群芽到 \(G'\) 的一个开李子群芽的 k-同构；并且 \(T_{e}(\phi)\) 是 K-线性的，因而当 x 充分接近 e 时 \(T_{x}(\phi)\) 是 K-线性的；所以 Int g 限制在 V 中 e 的某个充分小的开邻域上是 K-解析的（《可微与解析流形》，R，5.14.6）。根据§1第9号命题18，G 上存在一种解析 K-流形结构，使 G 成为一个李

K-群，且 V 是 G 的开子-K-流形。通过平移可见，G 的底层 k-流形结构就是给定的结构。李 K-群 G 的李代数与开李 K-子群芽 V 的李代数相同，因而为 L'。最后，推论中所述的唯一性由 § 3 第8号命题32推出。

定理 3. 设 G 是李群芽，\(\mathfrak{h}\) 是 \(\mathrm{L}(\mathrm{G})\) 中具有拓扑补空间的李子代数。存在 G 的李子群芽 H，使得 \(\mathrm{L}(\mathrm{H}) = \mathfrak{h}\)。若 \(\mathrm{H}_1\) 和 \(\mathrm{H}_2\) 是 G 的李子群芽，且

\[
\mathrm{L} (\mathrm{H} _ {1}) = \mathrm{L} (\mathrm{H} _ {2}) = \mathfrak {h},
\]

则 \(\mathrm{H}_1\cap \mathrm{H}_2\) 在 \(\mathrm{H}_1\) 和 \(\mathrm{H}_2\) 中都是开的。

存在一个李群芽 H'，其李代数与 \(\mathfrak{h}\) 同构（定理2）。必要时缩小 H'，可设存在从 H' 到 G 的态射 \(\phi\)，使得 \(\mathrm{L}(\phi)\) 是从 \(\mathrm{L}(\mathrm{H}')\) 到 \(\mathfrak{h}\) 的同构（第1号定理1）。由于 \(\mathfrak{h}\) 具有拓扑补空间，\(\phi\) 在 e 处是浸入。因此进一步缩小 H'，可设 \(\phi\) 是从流形 H' 到 G 的一个子流形的同构。这证明了 H 的存在性。第二个断言由下列命题推出：

命题 1. 设 G 是李群芽，H 和 H' 是两个李子群芽。要使 L(H) ⊃ L(H')，充要条件是 H ∩ H' 在 H' 中开。若 H ∩ H' 在 H' 中开，则 L(H') = L(H ∩ H') ⊂ L(H)。假设 L(H) ⊃ L(H')。设 i、i' 是 H、H' 到 G 的典范嵌入。必要时缩小 H'，可设存在从 H' 到 H 的态射 ψ，使得 L(ψ) 是 L(H') 到 L(H) 的典范嵌入（第1号定理1）。于是 L(i ∘ ψ) = L(i')，因此存在 H' 中 eH 的一个邻域 V，使得 i ∘ ψ 与 i' 在 V 上重合（定理1）。所以 V ⊂ H，进而 V ⊂ H ∩ H'，且 H ∩ H' 在 H' 中开（§ 1，第10号）。

命题 2. 设 G 是定义在 K 上的李群，k 是 K 的非离散闭子域，H 是李 k-群 G 的李子群。假设 L(H) 是 L(G) 的一个具有拓扑补空间的向量子-K-空间。则 H 是李 K-群 G 的李子群。

存在李 K-群 G 的一个李子群芽 \(\mathrm{H}'\)，使得 \(\mathrm{L}(\mathrm{H}') = \mathrm{L}(\mathrm{H})\)（定理3）。将 G、H、H' 看作李 \(k\)-群芽；定理3于是证明 \(\mathrm{H} \cap \mathrm{H}'\) 在 H 和 H' 中开。因此存在 G 中 e 的一个开邻域 U，使得 \(e\) 的 \(\mathrm{U} \cap \mathrm{H}\) 是 G 上的 K-子流形。所以 H 是李 K-群 G 的李子群（\(\S 1\)，第3号，命题6）。

